% Gemini theme
% https://github.com/anishathalye/gemini

\documentclass[final]{beamer}

% ====================
% Packages
% ====================

\usepackage[T1]{fontenc}
\usepackage{lmodern}
\usepackage[size=custom,width=84.1,height=118.9,scale=1.0]{beamerposter}
\usetheme{gemini}
\input{colorthemes/beamercolorthemegemini.sty}
\usepackage{graphicx}
\usepackage{booktabs}
\usepackage{array}
\usepackage{tikz}
\usepackage{pgfplots}
\pgfplotsset{compat=1.14}
\usepackage{anyfontsize}
\usepackage{amsmath}
\usepackage{bm}

% ====================
% Colors
% ====================

\definecolor{UniboRed}{RGB}{187, 46, 41}
\definecolor{UniboDark}{RGB}{92, 24, 25}
\definecolor{UniboRose}{RGB}{253, 242, 241}
\definecolor{UniboInk}{RGB}{35, 31, 32}
\definecolor{UniboMuted}{RGB}{97, 83, 82}
\definecolor{UniboGold}{RGB}{207, 166, 92}

\setbeamercolor{headline}{fg=white,bg=UniboRed}
\setbeamercolor{headline rule}{bg=UniboDark}
\setbeamercolor{structure}{fg=UniboRed}
\setbeamercolor{block title}{fg=UniboRed,bg=white}
\setbeamercolor{block separator}{bg=UniboRed}
\setbeamercolor{block body}{fg=UniboInk,bg=white}
\setbeamercolor{block title alerted}{fg=white,bg=UniboRed}
\setbeamercolor{block separator alerted}{bg=UniboDark}
\setbeamercolor{block body alerted}{fg=UniboInk,bg=UniboRose}
\setbeamercolor{block title example}{fg=UniboDark,bg=UniboRose}
\setbeamercolor{block separator example}{bg=UniboGold}
\setbeamercolor{block body example}{fg=UniboInk,bg=UniboRose}
\setbeamercolor{heading}{fg=UniboDark}
\setbeamercolor{item}{fg=UniboRed}

% ====================
% Lengths and helpers
% ====================

\newlength{\sepwidth}
\newlength{\colwidth}
\setlength{\sepwidth}{0.03\paperwidth}
\setlength{\colwidth}{0.455\paperwidth}

\newcommand{\separatorcolumn}{\begin{column}{\sepwidth}\end{column}}
\newcommand{\paperdir}{../../my_papers/2025_acl___Supervised_Relation_Extraction_is_More_Efficient_When_Approached_as_Graph_Based_Dependency_Parsing/latex}
\newcommand{\figdir}{\paperdir/figs}
\newcommand{\assetdir}{../../.figs}
\newcommand{\metric}[3]{%
  \begin{tikzpicture}[baseline=(label.base)]
    \node[rounded corners=3pt, fill=#1, inner xsep=0.55em, inner ysep=0.32em] (label)
      {\strut\textcolor{white}{\bfseries #2} \textcolor{white}{#3}};
  \end{tikzpicture}}

% ====================
% Title
% ====================

\title{LLMs Underperform Graph-Based Parsers on Supervised Relation Extraction for Complex Graphs}

\author{Paolo Gajo \inst{1} \and Domenic Rosati \inst{2} \and Hassan Sajjad \inst{2} \and Alberto Barrón-Cedeño \inst{1}}

\institute[shortinst]{\inst{1} University of Bologna \samelineand \inst{2} Dalhousie University}

% ====================
% Footer and logos
% ====================

\footercontent{
  \large ACL 2026, San Diego, USA \hfill
  % Relation extraction as graph-based dependency parsing \hfill
  \href{mailto:paolo.gajo2@unibo.it}{paolo.gajo2@unibo.it}}

\logoleft{%
  % \makebox[24cm][l]{%
    \includegraphics[height=6cm]{\assetdir/logos/universities/Seal_of_the_University_of_Bologna_bg.pdf}%
    \hspace{0.5cm}%
    \includegraphics[height=6cm]{\assetdir/logos/universities/dal_shield.pdf}%
  % }%
  }
\logoright{
    \includegraphics[width=18cm]{\assetdir/logos/conferences/ACL2026_logo_v8.png}%
  }

% ====================
% Body
% ====================

\setbeamerfont{normal text}{size=\Large}
\setbeamerfont{normal text}{size=\Large}
\setbeamerfont{block body}{size=\Large}
\setbeamerfont{block body alerted}{size=\Large}
\setbeamerfont{block body example}{size=\Large}
\setbeamerfont{caption}{size=\Large}
\setbeamerfont{caption name}{size=\Large}

\setbeamerfont{block title}{size=\Large}
\setbeamerfont{block title alerted}{size=\Large}
\setbeamerfont{block title example}{size=\Large}

\begin{document}

\begin{frame}[t]
\begin{columns}[t]
\separatorcolumn

\begin{column}{\colwidth}

  \begin{alertblock}{Motivation}
    Relation extraction (RE) is fundamental for knowledge graph construction. While LLMs can carry out RE from in-context learning (ICL) examples, ICL performance is not satisfactory without domain fine-tuning. Here, we compare a state-of-the-art graph-based parser architecture to LLMs for RE in supervised settings.
  \end{alertblock}

  \begin{block}{Relation extraction as latent graph inference}
    RE can be framed as predicting labeled edges over a text graph.
    Each document is represented as RDF-style triples, i.e. $(\texttt{head},\ \texttt{relation},\ \texttt{dependent})$.
    Each triple can be seen as two nodes connected by an edge. Thus, the task can be framed as dependency parsing.
    % : infer nodes, labeled directed edges, and edge labels over the linguistic graph.

    \vspace{1cm}

    

    \begin{figure}
      \centering
      \includegraphics[width=\linewidth]{\figdir/parser_diagram.pdf}
      \vspace{0.5cm}
      \caption{The graph-based parser scores head--dependent edges and relation labels with biaffine attention.}
      \label{fig:graph-parser-diagram}
    \end{figure}

    % \vspace{1.5cm}

    The graph parser we use scores edges via biaffine attention. Conversely, LLMs output their predictions by generating a JSON containing literal triples.

  \end{block}

  \begin{exampleblock}{Experimental setup}
    
    \vspace{0.4em}
    \centering
    \metric{UniboRed}{124\,M}{graph parser}
    \hspace{1em}
    \metric{UniboRed}{7\,B--70\,B}{LLMs}
    \hspace{1em}
    \metric{UniboRed}{6}{datasets}

    \vspace{0.4em}
    \raggedright

    \begin{itemize}
      \item \textbf{Graph parser:} frozen BERT with LSTMs and biaffine attn. (Fig.~\ref{fig:graph-parser-diagram}).
      \item \textbf{LLMs:} Mistral-7B-Instruct, Qwen3-14B-Base, Qwen3-32B, and Llama-3.3-70B-Instruct.
      \item \textbf{LLM Prompts:} Descriptive, no task info, adversarial.
    \end{itemize}
  \end{exampleblock}

  \begin{block}{Graph complexity varies widely}
    The datasets range from short, sparse relation graphs to recipe and dependency graphs with a high number $k$ of edges per document.

    \begin{table}
      \centering
      % \renewcommand{\arraystretch}{1.12}
      \vspace{0.5cm}
      \begin{tabular}{l@{\hspace{20mm}}r@{\hspace{20mm}}r@{\hspace{20mm}}r@{\hspace{20mm}}r@{\hspace{20mm}}}
        \toprule
        \textbf{Dataset} & \textbf{Min} & \textbf{Mean $k$} & \textbf{Max} & \textbf{$k \leq 5$} \\
        \midrule
        CoNLL04 & 1 & 1.42 & 11 & 98.54\% \\
        ADE     & 1 & 1.59 & 21 & 97.77\% \\
        SciERC  & 1 & 2.38 & 17 & 93.49\% \\
        enEWT   & 5 & 17.83 & 158 & 25.03\% \\
        SciDTB  & 5 & 23.41 & 100 & 2.50\% \\
        ERFGC   & 4 & 49.19 & 167 & 3.33\% \\
        \bottomrule
      \end{tabular}
      \vspace{0.5cm}
      \caption{Relation-count statistics across train, development, and test splits.}
    \end{table}

    % The core comparison asks whether model performance degrades as $k$ increases.
  \end{block}

\end{column}

\separatorcolumn

\begin{column}{\colwidth}

  \begin{block}{Results}
    After fine-tuning, LLMs perform well on small relation graphs, but their F$_1$ drops as graph complexity grows.
    The graph-based parser is lighter, faster at inference, and increasingly stronger on complex graphs.
  
    \begin{table}
      \centering
      % \renewcommand{\arraystretch}{1.14}
      \vspace{0.5cm}
      \begin{tabular}{l@{\hspace{10mm}}c@{\hspace{10mm}}c@{\hspace{10mm}}c@{\hspace{10mm}}c@{\hspace{10mm}}c@{\hspace{10mm}}c@{\hspace{10mm}}}
        \toprule
        \textbf{Model} & \textbf{CoNLL} & \textbf{ADE} & \textbf{SciERC} & \textbf{enEWT} & \textbf{SciDTB} & \textbf{ERFGC} \\
        \midrule
        Best LLM & \textbf{0.717} & \textbf{0.836} & \textbf{0.444} & 0.851 & 0.886 & 0.606 \\
        Graph parser & 0.668 & 0.697 & 0.351 & \textbf{0.865} & \textbf{0.918} & \textbf{0.713} \\
        % \midrule
        % Parser gain & -0.049 & -0.139 & -0.093 & +0.014 & +0.032 & +0.107 \\
        \bottomrule
      \end{tabular}
      \vspace{0.5cm}
      \caption{Best graph-based parser vs LLM results in terms of micro-F$_1$.}
    \end{table}

  \end{block}

  \begin{alertblock}{Graph-based parser vs best LLM}

    \vspace{0.5cm}

    \begin{tikzpicture}
      \centering
      \begin{axis}[
        width=0.9\linewidth,
        height=0.25\linewidth,
        font=\large,
        ybar,
        bar width=20pt,
        ymin=-0.25,
        ymax=0.25,
        ylabel={$\Delta$F$_1$},
        ylabel style={font=\Large},
        ytick={-0.15,0,0.15},
        symbolic x coords={CoNLL04,ADE,SciERC,enEWT,SciDTB,ERFGC},
        xtick=data,
        x tick label style={rotate=0, anchor=center, yshift=-0.5cm},
        y tick label style={rotate=0, anchor=east, font=\large, xshift=0cm},
        ymajorgrids=true,
        grid style={UniboRose},
        axis line style={UniboMuted},
        tick style={UniboMuted},
        nodes near coords,
        nodes near coords style={font=\large, /pgf/number format/fixed, /pgf/number format/precision=3},
      ]
        \addplot[fill=UniboRed, draw=UniboDark] coordinates {
          (CoNLL04,-0.049) (ADE,-0.139) (SciERC,-0.093)
          (enEWT,0.014) (SciDTB,0.032) (ERFGC,0.107)
        };
      \draw[dashed, thick, black]
        (rel axis cs:0,0.5) -- (rel axis cs:1,0.5);
      \end{axis}
    \end{tikzpicture}
  
  LLMs win on the sparse RE datasets. The parser wins once the average document has roughly 18 or more relations.
  
  \end{alertblock}

  \begin{block}{LLMs degrade more as graphs grow}
    The strongest LLM setting, Qwen3-14B-Base with a no-task-info prompt, shows a negative correlation between edge count $k$ and F$_1$ on five of six datasets. The parser is less sensitive to $k$.

    \begin{table}
      \centering
      % \renewcommand{\arraystretch}{1.14}
      \vspace{0.5cm}
      \begin{tabular}{l@{\hspace{10mm}}r@{\hspace{10mm}}r@{\hspace{10mm}}r@{\hspace{10mm}}r@{\hspace{10mm}}r@{\hspace{10mm}}r@{\hspace{10mm}}}
        \toprule
        \textbf{Model} & \textbf{CoNLL} & \textbf{ADE} & \textbf{SciERC} & \textbf{enEWT} & \textbf{SciDTB} & \textbf{ERFGC} \\
        \midrule
        Qwen3-14B & -0.143 & -0.152 & 0.171 & -0.270 & -0.185 & -0.639 \\
        Graph parser & 0.012 & -0.046 & 0.067 & -0.043 & -0.195 & -0.206 \\
        \bottomrule
      \end{tabular}
      \vspace{0.5cm}
      \caption{Pearson $r$ between relation count $k$ and micro-F$_1$.}
    \end{table}

    On ERFGC, the LLM performance is strongly anti-correlated with graph size, while the parser remains much more stable.
  \end{block}

  \begin{block}{Discussion}
    \begin{itemize}
      \item \textbf{LLMs serialize graph prediction.} Inputs, labels, and output formatting are all represented as text, so relevant entities and relation labels can become far apart in the prompt and generation.
      \item \textbf{Graph parsers score edges directly.} Pairwise scoring
      % keeps the task in the graph domain and
      avoids generating thousands of output tokens for dense graphs.
      \item \textbf{Fine-tuning overrides prompt info.} The no-task-info prompt layou is competitive with descriptive prompts.
      \item \textbf{Scale is not enough.} The 70\,B LLM improves on some datasets with more steps, but still trails the 124\,M parser on the most complex graphs.
    \end{itemize}
  \end{block}

  \begin{exampleblock}{Takeaways}
    \heading{For supervised relation extraction}
    Use graph-based parsers when documents contain many relations, especially dependency-like or recipe-flow graphs.

    \heading{For LLM-based extraction}
    The bottleneck is not only model capacity. It is also the mismatch between graph-structured supervision and text-serialized prediction.

    \heading{Future work}
    Investigate attention-level failure modes and reduce formatting overhead through compressed prompts or graph-native decoding.
  \end{exampleblock}

  % \begin{block}{Key references}
  %   \small
  %   Dozat and Manning (2017), Deep biaffine attention for neural dependency parsing. \\
  %   Bhatt et al. (2024), end-to-end graph-based parsing for entity and relation extraction. \\
  %   Gajo et al. (2026), LLMs underperform graph-based parsers on supervised relation extraction for complex graphs.
  % \end{block}

\end{column}

\separatorcolumn
\end{columns}
\end{frame}

\end{document}
